% ============================================================
%  ANEXOS
%  Cada anexo tiene un solo título (sin negrilla) y se desarrolla
%  en prosa continua, sin subtítulos internos. Las tablas y los
%  códigos se numeran por anexo: Tabla A.1, Código B.2, etc.
% ============================================================

\section*{\normalfont ANEXO I. Esquema completo del archivo manifest.json}
\addcontentsline{toc}{section}{ANEXO I. Esquema completo del archivo manifest.json}
\label{anx:manifest_schema}

\renewcommand{\thetable}{A.\arabic{table}}
\setcounter{table}{0}
\renewcommand{\thecodigo}{A.\arabic{codigo}}
\setcounter{codigo}{0}
\renewcommand{\thefigure}{A.\arabic{figure}}
\setcounter{figure}{0}

El archivo \textbf{manifest.json} funciona como el contrato de empaquetado de cada modelo. Debe ubicarse en la raíz del paquete ZIP para que el middleware pueda reconocerlo, validar su estructura e instalarlo de forma controlada. Durante la instalación, el servicio registra su contenido en \textbf{data/models\_store/registry.json} y guarda una copia en \textbf{data/bootstrap\_manifests/}; de esta manera, el registro puede reconstruirse automáticamente si se pierde o se reinicia el entorno.

El documento JSON se organiza en campos de primer nivel que describen la identidad del modelo, su forma de ejecución, sus artefactos y los valores que verá el usuario en ELAN: \textbf{model\_id} (identificador único dentro del registro), \textbf{name} (nombre legible mostrado en ELAN), \textbf{version} (versión semántica, por ejemplo \textbf{1.0.0}), \textbf{task} (tipo de tarea que realiza el modelo), \textbf{runtime} (modo de ejecución y tipo de runner), \textbf{artifacts} (mapa de rutas relativas a los artefactos del paquete), \textbf{input\_contract} y \textbf{output\_contract} (medio de entrada y tipo de salida), \textbf{container} (configuración de red del contenedor, obligatoria cuando se usa el runner \textbf{docker\_http}) y \textbf{ui} (valores predeterminados para el panel de configuración de ELAN). La Tabla~\ref{tab:manifest_root} detalla cada campo del objeto raíz, su tipo y su obligatoriedad.

\begin{table}[H]
\centering
\caption{Campos del objeto raíz del \texttt{manifest.json}}
\label{tab:manifest_root}
\scriptsize
\renewcommand{\arraystretch}{1.08}
\setlength{\tabcolsep}{3.5pt}

\begin{tabular}{
    >{\raggedright\arraybackslash}p{3.0cm}
    >{\raggedright\arraybackslash}p{1.6cm}
    >{\centering\arraybackslash}p{1.8cm}
    >{\raggedright\arraybackslash}p{6.3cm}
}
\toprule
\textbf{Campo} & \textbf{Tipo} & \textbf{Requerido} & \textbf{Descripción} \\
\midrule
\texttt{model\_id} & string & Sí & Identificador único; admite letras, dígitos, punto, guion y guion bajo (\texttt{[A-Za-z0-9\_.-]}). \\
\texttt{name} & string & Sí & Nombre descriptivo mostrado en la interfaz de ELAN. \\
\texttt{version} & string & Sí & Versión semántica del modelo, por ejemplo, \texttt{X.Y.Z}. \\
\texttt{task} & string & Sí & Tipo de tarea que resuelve el modelo. \\
\texttt{runtime} & object & Sí & Configuración del modo de ejecución. \\
\texttt{artifacts} & object & Sí & Mapa de rutas relativas a los artefactos del paquete. \\
\texttt{input\_contract} & object & Sí & Restricciones sobre el medio de entrada (tipo de medio, características). \\
\texttt{output\_contract} & object & Sí & Tipo de salida y clases esperadas. \\
\texttt{backend\_config} & object & No & Datos usados durante la instalación: nombre y tag de la imagen Docker, puerto y timeout de arranque inicial (180~s en el modelo de referencia). \\
\texttt{container} & object & Cond. & Obligatorio si \texttt{runtime.runner = docker\_http}. \\
\texttt{ui} & object & Sí & Valores predeterminados para el panel de ELAN; sus campos internos son opcionales. \\
\bottomrule
\end{tabular}
\end{table}

Respecto del campo \textbf{task}, el esquema \textbf{ModelManifest} del middleware acepta exactamente dos valores; cualquier otro provoca el rechazo del paquete durante la validación. El valor \textbf{video\_segmentation} indica que el modelo detecta regiones temporales en el video sin clasificar su contenido semántico, mientras que \textbf{video\_segmentation\_and\_gloss\_classification} indica que, además de detectar las regiones, el modelo asigna una glosa de lengua de señas a cada segmento; este último valor corresponde al modelo de referencia del trabajo.

El objeto \textbf{runtime} declara cómo debe ejecutarse el backend. Sus tres campos se resumen en la Tabla~\ref{tab:manifest_runtime}: el modo de ejecución (actualmente solo se soporta \textbf{docker}), el tipo de contenedor (reservado para extensiones futuras) y el identificador del runner del middleware encargado de iniciar el contenedor y llamar a los endpoints del backend.

\begin{table}[H]
\centering
\caption{Campos del objeto \texttt{runtime}}
\label{tab:manifest_runtime}
\scriptsize
\renewcommand{\arraystretch}{1.08}
\setlength{\tabcolsep}{3.5pt}

\begin{tabular}{
    >{\raggedright\arraybackslash}p{2.5cm}
    >{\raggedright\arraybackslash}p{1.7cm}
    >{\raggedright\arraybackslash}p{2.7cm}
    >{\raggedright\arraybackslash}p{6.5cm}
}
\toprule
\textbf{Campo} & \textbf{Tipo} & \textbf{Valor} & \textbf{Descripción} \\
\midrule
\texttt{mode} & string & \texttt{docker} & Modo de ejecución del backend. Actualmente, solo se soporta \texttt{docker}. \\
\texttt{framework} & string & \texttt{container} & Tipo de contenedor. Se mantiene reservado para extensiones futuras. \\
\texttt{runner} & string & \texttt{docker\_http} & Identifica el runner del middleware encargado de iniciar el contenedor y llamar a los endpoints \texttt{/health} e \texttt{/infer}. \\
\bottomrule
\end{tabular}
\end{table}

En cuanto al objeto \textbf{artifacts}, el middleware no interpreta el contenido de los artefactos; su función es conservar las rutas declaradas para que el \textbf{Dockerfile} del paquete pueda construir la imagen del backend con los archivos correctos. Las claves del mapa actúan como nombres semánticos definidos por el autor del modelo, mientras que los valores apuntan a rutas relativas dentro del paquete. Durante la instalación se comprueba que cada archivo declarado exista dentro del ZIP y que ninguna ruta sea absoluta ni contenga \textbf{../}; si alguna condición falla, el paquete se rechaza antes de construir la imagen. La Tabla~\ref{tab:manifest_artifacts} enumera los artefactos del modelo de referencia y su propósito.

\begin{table}[H]
\centering
\caption{Artefactos del modelo de referencia \texttt{lsec\_bio\_gloss\_final\_v1}}
\label{tab:manifest_artifacts}
\scriptsize
\renewcommand{\arraystretch}{1.08}
\setlength{\tabcolsep}{4pt}

\begin{tabular}{
    >{\raggedright\arraybackslash}p{3.1cm}
    >{\raggedright\arraybackslash}p{9.8cm}
}
\toprule
\textbf{Clave} & \textbf{Ruta relativa y propósito} \\
\midrule
\texttt{weights\_bio} & \texttt{weights/best\_bio\_segmenter\_v2.pt}: pesos del modelo BiLSTM de segmentación BIO. \\
\texttt{weights\_gloss} & \texttt{weights/best\_keypoint\_transformer\_v11.pt}: pesos del clasificador Transformer. \\
\texttt{vocab} & \texttt{vocab/gloss\_vocab\_top20.csv}: vocabulario de 20 glosas LSEc. \\
\texttt{config} & \texttt{config/pipeline\_config.json}: hiperparámetros del pipeline. \\
\texttt{dockerfile} & \texttt{backend/Dockerfile}: instrucciones de construcción de la imagen Docker. \\
\texttt{requirements} & \texttt{backend/requirements.txt}: dependencias Python del backend. \\
\texttt{app\_main} & \texttt{backend/app/main.py}: punto de entrada FastAPI del servicio de inferencia. \\
\bottomrule
\end{tabular}
\end{table}

El objeto \textbf{container} es obligatorio cuando \textbf{runtime.runner = docker\_http}, porque reúne los datos que el \textbf{DockerRunner} necesita para iniciar el backend, comprobar su disponibilidad y enviarle las solicitudes de inferencia. Sus campos se describen en la Tabla~\ref{tab:manifest_container}.

\begin{table}[H]
\centering
\caption{Campos del objeto \texttt{container}}
\label{tab:manifest_container}
\scriptsize
\renewcommand{\arraystretch}{1.08}
\setlength{\tabcolsep}{3.5pt}

\begin{tabular}{
    >{\raggedright\arraybackslash}p{3.3cm}
    >{\raggedright\arraybackslash}p{1.5cm}
    >{\raggedright\arraybackslash}p{2.4cm}
    >{\raggedright\arraybackslash}p{6.0cm}
}
\toprule
\textbf{Campo} & \textbf{Tipo} & \textbf{Valor típico} & \textbf{Descripción} \\
\midrule
\texttt{internal\_port} & integer & \texttt{8080} & Puerto expuesto por el backend dentro del contenedor. \\
\texttt{host\_port} & integer & \textit{opcional} & Puerto en el host para depuración; se recomienda omitirlo en producción. \\
\texttt{health\_path} & string & \texttt{/health} & Ruta del endpoint de comprobación de disponibilidad. \\
\texttt{infer\_path} & string & \texttt{/infer} & Ruta del endpoint de inferencia. \\
\texttt{startup\_timeout\_sec} & integer & \texttt{30} & Tiempo máximo, en segundos, que el middleware espera hasta que \texttt{/health} responda con estado 200 después de iniciar el contenedor. \\
\texttt{service\_url} & string & \textit{opcional} & URL base del servicio cuando el contenedor se ejecuta fuera del control del SDK. Si está presente, el runner omite la gestión del ciclo de vida. \\
\texttt{name} & string & \textit{opcional} & Nombre fijo del contenedor. Si se omite, Docker asigna un nombre automáticamente. \\
\texttt{environment} & object & \textit{opcional} & Variables de entorno pasadas al contenedor en tiempo de ejecución. \\
\texttt{volumes} & object & \textit{opcional} & Bind mounts adicionales del modelo. El middleware agrega automáticamente el volumen de videos. \\
\bottomrule
\end{tabular}
\end{table}

Por último, los valores del objeto \textbf{ui} se cargan automáticamente en el panel \textbf{AIOrchestrationRecognizerPanel} cuando el usuario selecciona un modelo en ELAN. Con ello se evita repetir configuraciones básicas en cada ejecución y se mantiene una experiencia más cercana al flujo habitual de anotación. Sus campos se resumen en la Tabla~\ref{tab:manifest_ui}.

\begin{table}[H]
\centering
\caption{Campos del objeto \texttt{ui}}
\label{tab:manifest_ui}
\scriptsize
\renewcommand{\arraystretch}{1.08}
\setlength{\tabcolsep}{4pt}

\begin{tabular}{
    >{\raggedright\arraybackslash}p{3.5cm}
    >{\raggedright\arraybackslash}p{1.7cm}
    >{\raggedright\arraybackslash}p{8.0cm}
}
\toprule
\textbf{Campo} & \textbf{Tipo} & \textbf{Descripción} \\
\midrule
\texttt{default\_label} & string & Etiqueta usada cuando \texttt{label\_mode = simple}. También actúa como valor de respaldo si el modelo no devuelve una glosa. \\
\texttt{default\_target\_tier} & string & Nombre del nivel de anotación de ELAN donde se insertarán las anotaciones. \\
\texttt{label\_mode} & string & Modo de generación de etiquetas. El valor \texttt{gloss\_top1} usa la glosa de mayor probabilidad, mientras que \texttt{simple} usa siempre \texttt{default\_label}. \\
\texttt{supports\_threshold} & boolean & Indica si el panel muestra un control de umbral de confianza. Actualmente, no es utilizado por el backend de referencia. \\
\bottomrule
\end{tabular}
\end{table}

Para cerrar el anexo, el Código~\ref{fig:manifest_ejemplo} presenta el \textbf{manifest.json} completo del modelo de referencia utilizado en la validación del sistema. Al incluirse en la raíz del paquete ZIP, este archivo se convierte en la fuente de información que guía la instalación, el registro del modelo y la configuración inicial mostrada en ELAN. Nótese la diferencia entre los dos timeouts de arranque: \textbf{backend\_config} usa 180 segundos porque durante la instalación el backend carga los pesos por primera vez, mientras que \textbf{container} usa 30 segundos porque en cada inferencia el contenedor ya debería estar preparado.

\begin{center}
\begin{minipage}{0.96\textwidth}
\hrule
\vspace{0.06cm}
\scriptsize
\begin{Verbatim}[baselinestretch=0.68]
{
  "model_id": "lsec_bio_gloss_final_v1",
  "name": "LSEC BIO Gloss Pipeline - Implementacion Final Tesis",
  "version": "1.0.0",
  "task": "video_segmentation_and_gloss_classification",
  "runtime": {
    "mode": "docker",
    "framework": "container",
    "runner": "docker_http"
  },
  "artifacts": {
    "weights_bio": "weights/best_bio_segmenter_v2.pt",
    "weights_gloss": "weights/best_keypoint_transformer_v11.pt",
    "vocab": "vocab/gloss_vocab_top20.csv",
    "config": "config/pipeline_config.json",
    "dockerfile": "backend/Dockerfile",
    "requirements": "backend/requirements.txt",
    "app_main": "backend/app/main.py"
  },
  "input_contract": {
    "media_type": "video",
    "feature_type": "keypoints",
    "input_dim": 356
  },
  "output_contract": {
    "type": "segments_with_gloss",
    "classes": ["O", "B", "I"],
    "top_k": 5
  },
  "ui": {
    "default_label": "LSEC_REGION",
    "default_target_tier": "AUTO_GLOSS_SEGMENTS",
    "label_mode": "gloss_top1",
    "supports_threshold": false
  },
  "backend_config": {
    "docker_image_name": "lsec-bio-gloss-final",
    "docker_image_tag": "1.0.0",
    "container_port": 8080,
    "health_path": "/health",
    "infer_path": "/infer",
    "startup_timeout_sec": 180
  },
  "container": {
    "internal_port": 8080,
    "health_path": "/health",
    "infer_path": "/infer",
    "startup_timeout_sec": 30
  }
}
\end{Verbatim}
\vspace{0.03cm}
\hrule

\captionof{codigo}{Manifest del modelo \texttt{lsec\_bio\_gloss\_final\_v1} v1.0.0}
\label{fig:manifest_ejemplo}
\end{minipage}
\end{center}

% ============================================================
%  Anexo II — Ejemplo completo de respuesta JSON de inferencia
% ============================================================

\section*{\normalfont ANEXO II. Ejemplo completo de respuesta JSON de inferencia}
\addcontentsline{toc}{section}{ANEXO II. Ejemplo completo de respuesta JSON de inferencia}
\label{anx:respuesta_json}

\renewcommand{\thetable}{B.\arabic{table}}
\setcounter{table}{0}
\renewcommand{\thecodigo}{B.\arabic{codigo}}
\setcounter{codigo}{0}
\renewcommand{\thefigure}{B.\arabic{figure}}
\setcounter{figure}{0}

Este anexo muestra una respuesta exitosa del endpoint \textbf{POST /api/v1/jobs/segment-video} del middleware. El ejemplo conserva los campos definidos en el esquema \textbf{SegmentVideoResponse} del archivo \textbf{middleware/app/schemas/jobs.py} y corresponde a una inferencia del modelo \textbf{lsec\_bio\_gloss\_final\_v1} v1.0.0 sobre un video corto de prueba de la LSEc. La intención es dejar claro cómo viajan los resultados desde el backend de IA hasta el formato que luego puede transformarse en anotaciones de ELAN.

El intercambio comienza con la solicitud que el cliente envía al middleware, reproducida en el Código~\ref{fig:request_inferencia}.

\begin{center}
\begin{minipage}{0.96\textwidth}
\hrule
\vspace{0.06cm}
\scriptsize
\begin{Verbatim}[baselinestretch=0.68]
POST http://127.0.0.1:8000/api/v1/jobs/segment-video
Content-Type: application/json

{
  "job_id": "job-elan-20260515143022512",
  "media": {
    "path": "/data/videos/lsec_validacion_001.mp4"
  },
  "annotation": {
    "target_tier": "AUTO_GLOSS_SEGMENTS",
    "default_label": "LSEC_REGION",
    "label_mode": "gloss_top1"
  },
  "model": {
    "model_id": "lsec_bio_gloss_final_v1",
    "version": "1.0.0"
  },
  "execution": {
    "device_preference": "auto",
    "runner": "auto",
    "timeout_sec": 300
  }
}
\end{Verbatim}
\vspace{0.03cm}
\hrule

\captionof{codigo}{Request de inferencia enviado al endpoint \texttt{POST /api/v1/jobs/segment-video}}
\label{fig:request_inferencia}
\end{minipage}
\end{center}

La respuesta correspondiente, con estado HTTP 200, incluye cinco campos de primer nivel: \textbf{job\_id}, \textbf{status}, \textbf{media\_info}, \textbf{segments} y \textbf{trace}. Los cuatro primeros identifican el trabajo, su estado, el video procesado y los segmentos obtenidos. El campo \textbf{trace}, por su parte, agrega el historial de estados, los tiempos de cada etapa del pipeline y los metadatos necesarios para observar la ejecución del contenedor Docker. Por brevedad, el Código~\ref{fig:respuesta_inferencia} muestra solo el primero de los cuatro segmentos detectados (\textbf{n\_detected\_segments}); los restantes siguen exactamente la misma estructura.

\begin{center}
\begin{minipage}{0.9\textwidth}
\hrule
\vspace{0.09cm}
\tiny
\begin{Verbatim}[baselinestretch=0.92]
HTTP/1.1 200 OK
Content-Type: application/json

{
  "job_id": "job-elan-20260515143022512",
  "status": "COMPLETED",

  "media_info": {
    "fps": 25.0,
    "duration_ms": 9840,
    "total_frames": 246
  },

  "segments": [
    {
      "segment_id": 1,
      "start_ms": 480,
      "end_ms": 1320,
      "start_frame": 12,
      "end_frame": 33,
      "duration_frames": 21,
      "label": "YO",
      "confidence": 0.9231,
      "predictions": [
        {"rank": 1, "gloss_id": 0,  "gloss": "YO",      "probability": 0.9231},
        {"rank": 2, "gloss_id": 6,  "gloss": "PERSONA", "probability": 0.0412},
        {"rank": 3, "gloss_id": 12, "gloss": "HOMBRE",  "probability": 0.0189},
        {"rank": 4, "gloss_id": 2,  "gloss": "QUERER",  "probability": 0.0091},
        {"rank": 5, "gloss_id": 9,  "gloss": "PENSAR",  "probability": 0.0077}
      ]
    }
  ],
  "trace": {
    "runner": "docker_http",
    "device": "container",
    "model_id": "lsec_bio_gloss_final_v1",
    "model_version": "1.0.0",
    "output_type": "segments_with_gloss",
    "exec_ms": 18472,
    "stages": {
      "validation_ms": 12,
      "queue_ms": 0,
      "container_start_ms": 3210,
      "healthcheck_ms": 842,
      "docker_inference_ms": 14356,
      "inference_ms": 14356,
      "keypoint_extraction_ms": 7821,
      "bio_model_load_ms": 390,
      "gloss_model_load_ms": 218,
      "vocab_load_ms": 44,
      "bio_inference_ms": 4103,
      "bio_postprocessing_ms": 312,
      "gloss_classification_ms": 1468,
      "postprocessing_ms": 52,
      "total_ms": 18472
    },
    "state_history": [
      "RECEIVED",
      "VALIDATING",
      "PREPROCESSING",
      "QUEUED",
      "RUNNING",
      "POSTPROCESSING",
      "COMPLETED"
    ],
    "fps": 25.0,
    "total_frames": 246,
    "sampled_frames": 246,
    "windows_count": 30,
    "original_width": 640,
    "original_height": 480,
    "n_detected_segments": 4,
    "docker_mode": "sdk_managed",
    "docker_image": "lsec-bio-gloss-final:1.0.0",
    "container_id": "a3f7b91c2d84e056f1a2b3c4d5e6f789abcd1234ef015678cd9012ab34d63456",
    "container_start_ms": 3210,
    "healthcheck_ms": 842,
    "docker_inference_ms": 14356,
    "total_ms": 18472
  }
}
\end{Verbatim}
\vspace{0.03cm}
\hrule

\captionof{codigo}{Respuesta completa de inferencia exitosa HTTP 200}
\label{fig:respuesta_inferencia}
\end{minipage}
\end{center}

Para interpretar la respuesta, la Tabla~\ref{tab:trace_campos} describe cada campo del objeto \textbf{trace}, que concentra la información de trazabilidad de la ejecución.

\begin{table}[H]
\centering
\caption{Campos del objeto \texttt{trace} en la respuesta de inferencia}
\label{tab:trace_campos}
\scriptsize
\renewcommand{\arraystretch}{1.08}
\setlength{\tabcolsep}{4pt}

\begin{tabular}{
    >{\raggedright\arraybackslash}p{4.2cm}
    >{\raggedright\arraybackslash}p{8.8cm}
}
\toprule
\textbf{Campo} & \textbf{Descripción} \\
\midrule
\texttt{runner} & Identificador del runner del middleware, por ejemplo, \texttt{docker\_http}. \\
\texttt{device} & Tipo de dispositivo de ejecución reportado por el runner, por ejemplo, \texttt{container}. \\
\texttt{model\_id} & Identificador del modelo que procesó el job. \\
\texttt{model\_version} & Versión del modelo. \\
\texttt{output\_type} & Tipo de salida producida, por ejemplo, \texttt{segments\_with\_gloss}. \\
\texttt{exec\_ms} & Tiempo total de ejecución en el middleware, expresado en milisegundos. \\
\texttt{stages} & Desglose de latencia por etapa del procesamiento. \\
\texttt{state\_history} & Lista ordenada de estados que atravesó el job. \\
\texttt{fps / total\_frames} & Metadatos del video procesado, propagados desde el backend. \\
\texttt{n\_detected\_segments} & Número de segmentos temporales detectados. \\
\texttt{docker\_mode} & Modo de gestión del contenedor, por ejemplo, \texttt{sdk\_managed} o \texttt{compose\_service}. \\
\texttt{docker\_image} & Imagen Docker utilizada por el backend de inferencia. \\
\texttt{container\_id} & Identificador del contenedor Docker. \\
\texttt{container\_start\_ms} & Tiempo empleado en iniciar el contenedor, expresado en milisegundos. \\
\texttt{healthcheck\_ms} & Tiempo hasta que \textbf{GET /health} respondió HTTP 200, expresado en milisegundos. \\
\texttt{docker\_inference\_ms} & Tiempo de la llamada \textbf{POST /infer} dentro del contenedor, expresado en milisegundos. \\
\bottomrule
\end{tabular}
\end{table}

Dentro de \textbf{trace}, el objeto \textbf{stages} corresponde al esquema \textbf{StageMetrics} definido en \textbf{middleware/app/schemas/metrics.py}. Su utilidad principal es separar la latencia total en dos niveles: las etapas propias del orquestador y las etapas internas del pipeline de IA. Esta separación facilita identificar si el tiempo de respuesta se consume en validación, espera, comunicación con Docker o procesamiento del modelo. El significado de cada campo se presenta en la Tabla~\ref{tab:stages_campos}.

\begin{table}[H]
\centering
\caption{Significado de los campos del objeto \texttt{stages}}
\label{tab:stages_campos}
\scriptsize
\renewcommand{\arraystretch}{1.08}
\setlength{\tabcolsep}{4pt}

\begin{tabular}{
    >{\raggedright\arraybackslash}p{4.8cm}
    >{\raggedright\arraybackslash}p{8.2cm}
}
\toprule
\textbf{Campo} & \textbf{Descripción} \\
\midrule
\texttt{validation\_ms} & Validación del cuerpo JSON del request mediante Pydantic. \\
\texttt{queue\_ms} & Tiempo de espera en la cola FIFO de jobs. \\
\texttt{container\_start\_ms} & Arranque del contenedor Docker mediante el SDK. \\
\texttt{healthcheck\_ms} & Espera activa hasta que \texttt{/health} responde 200. \\
\texttt{docker\_inference\_ms} & Tiempo total de la llamada \textbf{POST /infer} al backend. \\
\texttt{keypoint\_extraction\_ms} & Extracción de keypoints con MediaPipe Holistic. \\
\texttt{bio\_model\_load\_ms} & Carga del modelo BiLSTM a memoria. \\
\texttt{gloss\_model\_load\_ms} & Carga del Transformer de clasificación de glosas. \\
\texttt{vocab\_load\_ms} & Lectura del vocabulario CSV de glosas. \\
\texttt{bio\_inference\_ms} & Inferencia del segmentador BIO mediante ventana deslizante. \\
\texttt{bio\_postprocessing\_ms} & Suavizado y filtrado de etiquetas BIO. \\
\texttt{gloss\_classification\_ms} & Clasificación Transformer de cada segmento detectado. \\
\texttt{postprocessing\_ms} & Construcción de la respuesta final en el middleware. \\
\texttt{total\_ms} & Tiempo total de extremo a extremo. \\
\bottomrule
\end{tabular}
\end{table}

% ============================================================
%  Anexo III — Configuración completa del docker-compose.yml
% ============================================================

\section*{\normalfont ANEXO III. Configuración del docker-compose.yml del middleware}
\addcontentsline{toc}{section}{ANEXO III. Configuración del docker-compose.yml del middleware}
\label{anx:docker_compose}

\renewcommand{\thetable}{C.\arabic{table}}
\setcounter{table}{0}
\renewcommand{\thecodigo}{C.\arabic{codigo}}
\setcounter{codigo}{0}
\renewcommand{\thefigure}{C.\arabic{figure}}
\setcounter{figure}{0}

El archivo \textbf{docker-compose.yml} reúne la configuración necesaria para ejecutar el middleware como servicio local y conectarlo con los contenedores de modelos de IA. Es el mecanismo recomendado para el despliegue, porque concentra en un solo lugar el ciclo de vida del contenedor, las variables de entorno, los bind mounts de persistencia y la red compartida \textbf{elan-ai-shared}. El Código~\ref{fig:docker_compose} reproduce el archivo completo utilizado durante el desarrollo y la validación.

\begin{center}
\begin{minipage}{0.96\textwidth}
\hrule
\vspace{0.06cm}
\scriptsize
\begin{Verbatim}[baselinestretch=0.68]
services:
  middleware:
    build:
      context: .
    image: elan-ai-middleware:0.1.0
    container_name: elan-ai-middleware
    ports:
      - "127.0.0.1:8000:8000"
    environment:
      MIDDLEWARE_HOST: 0.0.0.0
      MIDDLEWARE_PORT: 8000
      MIDDLEWARE_LOG_LEVEL: INFO
      MIDDLEWARE_RUNTIME_PROFILE: final
      MIDDLEWARE_MODELS_STORE_DIR: /app/data/models_store
      # Bootstrap: on every startup the middleware scans this directory
      # and auto-registers any manifest.json found there.
      # Combined with the bind mount below, installed model manifests
      # survive even a full volume reset.
      MIDDLEWARE_BOOTSTRAP_MANIFESTS_DIR: /app/data/bootstrap_manifests
      # Directory on the HOST where videos are stored.
      # Mounted read-only at /data/videos inside every model container.
      # Use forward slashes even on Windows (Docker Desktop translates them).
      MIDDLEWARE_VIDEOS_DIR: "C:/Users/imbaq/OneDrive/Desktop"
      # Shared Docker network used for container-to-container communication.
      # Model containers are attached to this network so the middleware
      # can reach them by container name (no host-port mapping required).
      MIDDLEWARE_DOCKER_NETWORK: elan-ai-shared
    volumes:
      # BIND MOUNTS: live on the host filesystem and are NEVER removed
      # by "docker compose down -v", "docker volume prune" or Docker Desktop.
      - ./data/models_store:/app/data/models_store
      - ./data/bootstrap_manifests:/app/data/bootstrap_manifests
      # Mount the Docker socket so the middleware can build/run model containers.
      - /var/run/docker.sock:/var/run/docker.sock
    networks:
      - shared
    restart: unless-stopped

# No named volumes: everything is on the host filesystem via bind mounts.

networks:
  shared:
    # Fixed name so model containers can join by this exact name
    # regardless of the docker-compose project prefix.
    name: elan-ai-shared
\end{Verbatim}
\vspace{0.03cm}
\hrule

\captionof{codigo}{Archivo \texttt{docker-compose.yml} del middleware}
\label{fig:docker_compose}
\end{minipage}
\end{center}

Las variables de entorno del servicio se describen en la Tabla~\ref{tab:env_vars}. Las variables que no aparecen declaradas en el archivo, como \textbf{MIDDLEWARE\_MAX\_CONCURRENT\_JOBS}, toman su valor por defecto desde \textbf{app/core/config.py}; para modificarlas basta con añadirlas a la sección \textbf{environment} del servicio y reiniciar el contenedor.

\begin{table}[H]
\centering
\caption{Variables de entorno del servicio middleware}
\label{tab:env_vars}
\tiny
\renewcommand{\arraystretch}{1.08}
\setlength{\tabcolsep}{4pt}

\begin{tabular}{
    >{\raggedright\arraybackslash}p{4.6cm}
    >{\raggedright\arraybackslash}p{3.2cm}
    >{\raggedright\arraybackslash}p{5.6cm}
}
\toprule
\textbf{Variable} & \textbf{Valor por defecto} & \textbf{Descripción} \\
\midrule
\texttt{MIDDLEWARE\_HOST} &
\texttt{0.0.0.0} &
Interfaz en que Uvicorn acepta conexiones dentro del contenedor. \\

\texttt{MIDDLEWARE\_PORT} &
\texttt{8000} &
Puerto de escucha de Uvicorn. \\

\texttt{MIDDLEWARE\_LOG\_LEVEL} &
\texttt{INFO} &
Nivel de log de la aplicación FastAPI. \\

\texttt{MIDDLEWARE\_RUNTIME\_PROFILE} &
\texttt{final} &
Perfil de ejecución. En producción se utiliza \texttt{final}. \\

\texttt{MIDDLEWARE\_MODELS\_STORE\_DIR} &
\path{/app/data/models_store} &
Ruta al directorio del registro de modelos y archivos de instalación. \\

\texttt{MIDDLEWARE\_BOOTSTRAP\_MANIFESTS\_DIR} &
\path{/app/data/bootstrap_manifests} &
Directorio escaneado al iniciar para reconstruir el registro automáticamente. \\

\texttt{MIDDLEWARE\_VIDEOS\_DIR} &
\textit{requerido} &
Ruta del host donde residen los videos. Se monta como \texttt{/data/videos} en los contenedores de modelo. \\

\texttt{MIDDLEWARE\_DOCKER\_NETWORK} &
\texttt{elan-ai-shared} &
Nombre de la red Docker compartida. Los contenedores de modelo se conectan a esta red para ser alcanzables por su nombre. \\

\texttt{MIDDLEWARE\_MAX\_CONCURRENT\_JOBS} &
\texttt{1} &
Número máximo de inferencias simultáneas. Controla el acceso concurrente a GPU, CPU o memoria del host. \\
\bottomrule
\end{tabular}
\end{table}

La configuración utiliza únicamente bind mounts, es decir, directorios del equipo anfitrión montados dentro del contenedor. Esta decisión evita depender de volúmenes Docker con nombre y facilita revisar los archivos directamente desde el sistema operativo. Además, los datos de modelos y manifests se conservan incluso después de detener el servicio o limpiar volúmenes desde Docker. Los tres montajes utilizados y su propósito se detallan en la Tabla~\ref{tab:bind_mounts}.

\begin{table}[H]
\centering
\caption{Bind mounts del servicio middleware}
\label{tab:bind_mounts}
\tiny
\renewcommand{\arraystretch}{1.08}
\setlength{\tabcolsep}{3.5pt}

\begin{tabular}{
    >{\raggedright\arraybackslash}p{4.3cm}
    >{\raggedright\arraybackslash}p{4.3cm}
    >{\raggedright\arraybackslash}p{4.8cm}
}
\toprule
\textbf{Ruta host} & \textbf{Ruta contenedor} & \textbf{Propósito} \\
\midrule
\texttt{./data/models\_store} &
\texttt{/app/data/models\_store} &
Registro de modelos instalados, archivo \texttt{registry.json} y archivos del paquete de cada modelo. \\

\texttt{./data/bootstrap\_manifests} &
\path{/app/data/bootstrap_manifests} &
Copias de los manifests de modelos para reconstrucción automática del registro al reiniciar. \\

\texttt{/var/run/docker.sock} &
\texttt{/var/run/docker.sock} &
Socket del daemon Docker. Permite al middleware construir imágenes y gestionar contenedores de modelos mediante el SDK de Python. \\
\bottomrule
\end{tabular}
\end{table}

En cuanto a la red, \textbf{elan-ai-shared} conserva un nombre fijo, independiente del prefijo que Docker Compose asigna al proyecto. Esto permite que los contenedores de modelo, creados dinámicamente por el SDK de Docker durante la inferencia, se unan a la misma red sin configuraciones adicionales. Una vez conectados, el middleware puede comunicarse con cada backend por su nombre de contenedor mediante la resolución DNS interna de Docker, sin exponer puertos innecesarios en el host. Los comandos más frecuentes durante la operación del servicio se resumen en el Código~\ref{fig:compose_cmds}.

\begin{center}
\begin{minipage}{0.90\textwidth}
\hrule
\vspace{0.06cm}
\scriptsize
\begin{Verbatim}[baselinestretch=0.72]
# Construir la imagen e iniciar el middleware en segundo plano
docker compose up --build -d

# Ver logs en tiempo real
docker compose logs -f middleware

# Detener el middleware; los datos persisten en ./data/
docker compose down

# Detener y eliminar la imagen; los datos persisten en ./data/
docker compose down --rmi local

# Verificar que el middleware responde
curl http://127.0.0.1:8000/health

# Consultar modelos registrados
curl http://127.0.0.1:8000/api/v1/models

# Consultar metricas de operacion
curl http://127.0.0.1:8000/api/v1/metrics
\end{Verbatim}
\vspace{0.03cm}
\hrule

\captionof{codigo}{Comandos Docker Compose más frecuentes durante la operación}
\label{fig:compose_cmds}
\end{minipage}
\end{center}

Respecto de la seguridad y la red, el puerto 8000 del middleware se expone únicamente en la interfaz de loopback del host, mediante \textbf{127.0.0.1:8000}. Con esta decisión, el servicio queda disponible para ELAN en la misma máquina, pero no para otros equipos de la red local. Esta restricción encaja con el escenario previsto del TIC, en el que el investigador ejecuta ELAN, el middleware y Docker en su propio equipo de trabajo. Si se quisiera que otras estaciones con ELAN consumieran el mismo middleware, como se plantea en las recomendaciones, bastaría con cambiar la publicación del puerto de \textbf{"127.0.0.1:8000:8000"} a \textbf{"8000:8000"} y apuntar la URL del panel del reconocedor a la dirección IP del servidor; en ese caso se recomienda añadir mecanismos de autenticación y cifrado si la red no es de confianza, además de ajustar \textbf{MIDDLEWARE\_MAX\_CONCURRENT\_JOBS} a la capacidad del equipo.

El socket de Docker montado en el servicio entrega al contenedor del middleware un nivel de control amplio sobre el daemon Docker del host. Para un entorno local y controlado, esta decisión simplifica la construcción y ejecución de modelos; sin embargo, en escenarios compartidos o multiusuario conviene restringir ese acceso mediante un proxy de Docker, como \textbf{Tecnativa Docker Socket Proxy}, o migrar a la API REST de Docker Engine con autenticación TLS.

% ============================================================
%  Anexo IV — Instrucciones de compilación e integración en ELAN 7.1
% ============================================================

\section*{\normalfont ANEXO IV. Compilación e integración del reconocedor en ELAN 7.1}
\addcontentsline{toc}{section}{ANEXO IV. Compilación e integración del reconocedor en ELAN 7.1}
\label{anx:compilacion_elan}

\renewcommand{\thetable}{D.\arabic{table}}
\setcounter{table}{0}
\renewcommand{\thecodigo}{D.\arabic{codigo}}
\setcounter{codigo}{0}
\renewcommand{\thefigure}{D.\arabic{figure}}
\setcounter{figure}{0}

Este anexo describe cómo compilar el código fuente modificado de ELAN 7.1 e integrar el reconocedor de IA en una instalación funcional de la herramienta. La modificación incorpora el paquete \textbf{mpi.eudico.client.annotator.recognizer.ai} al árbol fuente de ELAN, manteniendo intactas las clases originales. Esta decisión reduce el riesgo de afectar funcionalidades existentes y facilita retirar o actualizar el reconocedor en el futuro.

Antes de compilar deben cumplirse los prerrequisitos de la Tabla~\ref{tab:prerrequisitos}. En particular, la variable de entorno \textbf{JAVA\_HOME} debe apuntar al JDK 21 antes de ejecutar Maven; esta verificación evita errores de compilación asociados a versiones de Java distintas de las declaradas en el proyecto.

\begin{table}[H]
\centering
\caption{Prerrequisitos para compilar ELAN 7.1 con el reconocedor de IA}
\label{tab:prerrequisitos}
\scriptsize
\renewcommand{\arraystretch}{1.08}
\setlength{\tabcolsep}{4pt}

\begin{tabular}{
    >{\raggedright\arraybackslash}p{3.3cm}
    >{\raggedright\arraybackslash}p{2.4cm}
    >{\raggedright\arraybackslash}p{7.4cm}
}
\toprule
\textbf{Herramienta} & \textbf{Versión mínima} & \textbf{Notas} \\
\midrule
JDK OpenJDK/Oracle &
21 &
El archivo \textbf{pom.xml} de ELAN especifica \textbf{source} y \textbf{target} en nivel 21. \\

Apache Maven &
3.8 &
Gestor de dependencias y ciclo de vida de construcción. \\

Git &
2.x &
Permite clonar o actualizar el repositorio fuente de ELAN. \\

Docker Desktop &
4.x &
Requerido en tiempo de ejecución para el middleware y los backends de modelos. \\
\bottomrule
\end{tabular}
\end{table}

Las clases Java añadidas se concentran en un solo paquete dentro del repositorio de ELAN. El Código~\ref{fig:paquete_java_ai} muestra la ubicación del reconocedor, su panel de configuración, el cliente HTTP y los DTO utilizados para intercambiar datos con el middleware.

\begin{center}
\begin{minipage}{0.96\textwidth}
\hrule
\vspace{0.06cm}
\scriptsize
\begin{Verbatim}[baselinestretch=0.68]
elan-7.1/src/main/java/mpi/eudico/client/annotator/recognizer/ai/
|
+-- AIOrchestrationRecognizer.java          [implementa Recognizer AVATecH]
+-- AIOrchestrationRecognizerPanel.java     [panel de configuracion Swing]
+-- AIOrchestrationMiddlewareClient.java    [cliente HTTP del middleware]
+-- AIOrchestrationMiddlewareException.java [excepcion tipada]
+-- AIModelManagerDialog.java               [dialogo de gestion de modelos]
|
+-- dto/
    +-- SegmentVideoRequest.java            [DTO de request de inferencia]
    +-- SegmentVideoResponse.java           [DTO de respuesta de inferencia]
    +-- Segment.java                        [segmento temporal]
    +-- Prediction.java                     [glosa candidata con probabilidad]
    +-- MediaInfo.java                      [metadatos del video]
    +-- Trace.java                          [informacion de traza del job]
    +-- ModelSummary.java                   [resumen de modelo instalado]
    +-- InstalledModelDetail.java           [detalle completo de modelo]
    +-- ModelUiConfig.java                  [configuracion UI del modelo]
    +-- MediaPayload.java                   [sub-DTO del campo media]
    +-- AnnotationPayload.java              [sub-DTO del campo annotation]
    +-- ModelPayload.java                   [sub-DTO del campo model]
    +-- ExecutionPayload.java               [sub-DTO del campo execution]
    +-- ParametersPayload.java              [sub-DTO del campo parameters]
\end{Verbatim}
\vspace{0.03cm}
\hrule

\captionof{codigo}{Paquete Java del reconocedor de IA dentro del árbol fuente de ELAN}
\label{fig:paquete_java_ai}
\end{minipage}
\end{center}

Sobre las dependencias, el cliente HTTP \textbf{AIOrchestrationMiddlewareClient} usa clases disponibles en el JDK 11+, en particular \textbf{java.net.http.HttpClient}, por lo que las llamadas al middleware no requieren librerías externas. Para el parseo de las respuestas JSON se utiliza \textbf{org.json}, dependencia que ya forma parte del archivo \textbf{pom.xml} de ELAN 7.1. En consecuencia, la integración no agrega nuevas dependencias al proyecto.

ELAN descubre los reconocedores disponibles mediante el mecanismo Service Provider Interface (SPI) de Java. Por ello, el nuevo reconocedor debe declararse en el archivo \textbf{src/main/resources/META-INF/services/mpi.eudico.client.annotator.recognizer.api.Recognizer} con la línea del Código~\ref{fig:spi_recognizer}, añadida al final del archivo existente sin eliminar otras implementaciones registradas. Al iniciar la función \textit{Run Recognizer}, ELAN lee este archivo y carga todos los reconocedores declarados.

\begin{center}
\begin{minipage}{0.92\textwidth}
\hrule
\vspace{0.06cm}
\scriptsize
\begin{Verbatim}[baselinestretch=0.72]
mpi.eudico.client.annotator.recognizer.ai.AIOrchestrationRecognizer
\end{Verbatim}
\vspace{0.03cm}
\hrule

\captionof{codigo}{Línea añadida al archivo SPI para registrar el reconocedor de IA}
\label{fig:spi_recognizer}
\end{minipage}
\end{center}

La compilación sigue el ciclo de vida estándar de Maven. Los comandos del Código~\ref{fig:maven_build} deben ejecutarse desde el directorio raíz del proyecto \textbf{elan-7.1/}, de modo que Maven encuentre el archivo \textbf{pom.xml}, resuelva las dependencias y genere el JAR final. Si la compilación termina correctamente, Maven reporta \textbf{BUILD SUCCESS} y deja el archivo generado en \textbf{target/elan-7.1.jar}.

\begin{center}
\begin{minipage}{0.88\textwidth}
\hrule
\vspace{0.06cm}
\scriptsize
\begin{Verbatim}[baselinestretch=0.72]
# 1. Verificar version de Java; debe ser 21
java -version

# 2. Limpiar compilaciones anteriores
mvn clean

# 3. Compilar el arbol fuente completo
mvn compile

# 4. Crear el JAR sin lanzar ELAN
mvn package -Dexec.phase=install

# El JAR resultante queda en:
#   elan-7.1/target/elan-7.1.jar

# Las dependencias se copian en:
#   elan-7.1/target/lib/
\end{Verbatim}
\vspace{0.03cm}
\hrule

\captionof{codigo}{Proceso de compilación y empaquetado con Maven}
\label{fig:maven_build}
\end{minipage}
\end{center}

Durante el desarrollo, ELAN puede ejecutarse directamente desde Maven con los comandos del Código~\ref{fig:maven_run}. Esta opción permite probar el reconocedor integrado sin reemplazar todavía el JAR de una instalación existente.

\begin{center}
\begin{minipage}{0.74\textwidth}
\hrule
\vspace{0.06cm}
\scriptsize
\begin{Verbatim}[baselinestretch=0.72]
# Compilar y ejecutar ELAN directamente
mvn test

# Ejecutar sin recompilar, si ya se compiló
mvn exec:exec
\end{Verbatim}
\vspace{0.03cm}
\hrule

\captionof{codigo}{Ejecución de ELAN durante el desarrollo sin construir el JAR}
\label{fig:maven_run}
\end{minipage}
\end{center}

Para distribuir ELAN con el reconocedor de IA integrado, se puede reemplazar el JAR original de una instalación existente siguiendo el Código~\ref{fig:reemplazar_jar}. Antes de hacerlo, conviene guardar una copia de seguridad para poder restaurar la herramienta si el nuevo paquete no inicia correctamente. Debe considerarse, además, que si la instalación de ELAN utiliza un \textit{launcher} independiente, como ocurre con algunos instaladores de MPI, el nombre del JAR puede variar; por esta razón, antes de reemplazarlo se debe revisar qué archivo ejecuta realmente el script de lanzamiento o el acceso directo de ELAN.

\begin{center}
\begin{minipage}{0.92\textwidth}
\hrule
\vspace{0.06cm}
\scriptsize
\begin{Verbatim}[baselinestretch=0.72]
# Ruta de instalacion tipica de ELAN en Windows
# C:\Program Files\ELAN_7-1\

# Hacer copia de seguridad del JAR original
copy "C:\Program Files\ELAN_7-1\ELAN.jar" `
     "C:\Program Files\ELAN_7-1\ELAN.jar.bak"

# Copiar el nuevo JAR compilado
copy "elan-7.1\target\elan-7.1.jar" `
     "C:\Program Files\ELAN_7-1\ELAN.jar"
\end{Verbatim}
\vspace{0.03cm}
\hrule

\captionof{codigo}{Sustitución del JAR en una instalación existente de ELAN en Windows}
\label{fig:reemplazar_jar}
\end{minipage}
\end{center}

Con ELAN ejecutándose desde el JAR modificado, la integración puede comprobarse siguiendo el flujo normal de uso del reconocedor:

\begin{enumerate}[noitemsep, topsep=2pt, parsep=0pt, partopsep=0pt]
    \item Abrir un archivo EAF en ELAN o crear uno nuevo.
    \item Navegar a \textbf{Recognizer} $\rightarrow$ \textbf{Run Recognizer\ldots}
    \item En el desplegable de reconocedores, verificar que aparece la entrada \textbf{AI Orchestration Middleware Bridge}.
    \item Seleccionar dicha entrada. El panel \textbf{AIOrchestrationRecognizerPanel} debe cargarse en la sección inferior del diálogo.
    \item Asegurarse de que el middleware está en ejecución mediante \textbf{docker compose up -d}.
    \item Pulsar el botón \textbf{Probar conexión} para verificar la conectividad con \textbf{GET /health}.
    \item Desde el botón \textbf{Gestionar modelos\ldots}, instalar el modelo \textbf{lsec\_bio\_gloss\_final\_v1} subiendo el paquete ZIP.
    \item Seleccionar el modelo instalado en el desplegable de modelos.
    \item Pulsar \textbf{Start} en el diálogo de reconocedores para iniciar la inferencia.
\end{enumerate}

Si la integración funciona correctamente, las anotaciones generadas por el modelo aparecerán al finalizar el proceso en un nuevo nivel de anotación llamado \textbf{AUTO\_GLOSS\_SEGMENTS}. Este resultado confirma que la inferencia se ejecutó fuera de ELAN, pero regresó a la herramienta como segmentos temporales editables por el investigador. Para finalizar, la Tabla~\ref{tab:troubleshooting} reúne los problemas más frecuentes observados durante la compilación y la integración, junto con su solución.

\begin{table}[H]
\centering
\caption{Problemas frecuentes durante la compilación e integración}
\label{tab:troubleshooting}
\scriptsize
\renewcommand{\arraystretch}{1.08}
\setlength{\tabcolsep}{4pt}

\begin{tabular}{
    >{\raggedright\arraybackslash}p{4.7cm}
    >{\raggedright\arraybackslash}p{8.4cm}
}
\toprule
\textbf{Problema} & \textbf{Solución} \\
\midrule
\textbf{JAVA\_HOME} no configurado &
Definir \textbf{JAVA\_HOME} apuntando al directorio raíz del JDK 21. \\

El reconocedor no aparece en ELAN &
Verificar que la línea SPI se añadió correctamente al archivo \textbf{META-INF/services/...Recognizer}. \\

\textbf{BUILD FAILURE: cannot find symbol} &
Verificar que todos los archivos del paquete \textbf{ai} y \textbf{ai/dto} están presentes en el árbol fuente. \\

ELAN no inicia tras reemplazar el JAR &
Restaurar \textbf{ELAN.jar.bak} y verificar que el nuevo JAR se compiló sin errores. \\

Panel muestra error de conexión &
Verificar que el middleware está en ejecución mediante \textbf{curl http://127.0.0.1:8000/health}. \\

Timeout durante la instalación del modelo &
El proceso de construcción de Docker puede tardar entre 5 y 10 minutos durante la primera instalación. El timeout del cliente está configurado en 300 segundos. \\

\textbf{422 Unprocessable Entity} al instalar &
El boundary de \textbf{multipart/form-data} no debe comenzar con \textbf{--}. Verificar la versión compilada de \textbf{AIOrchestrationMiddlewareClient.installModel()}. \\
\bottomrule
\end{tabular}
\end{table}

% ============================================================
%  Anexo E — Repositorio del código fuente del middleware
% ============================================================

\section*{\normalfont ANEXO E. Repositorio del código fuente del middleware}
\addcontentsline{toc}{section}{ANEXO E. Repositorio del código fuente del middleware}
\label{anx:repo_middleware}

El código fuente completo del middleware de orquestación desarrollado en este trabajo se encuentra publicado en el siguiente repositorio de GitHub:

% TODO: reemplazar por la URL real del repositorio
\begin{center}
\url{https://github.com/TPanchoX/middleware-orquestacion-elan}
\end{center}

El repositorio contiene la aplicación FastAPI organizada según la estructura de directorios presentada en la metodología (Código~\ref{lst:estructura_directorios}), junto con el \textbf{Dockerfile} del servicio, el archivo \textbf{docker-compose.yml} reproducido en el Anexo~III y las dependencias de Python declaradas en \textbf{requirements.txt}. Con este material, cualquier lector puede clonar el proyecto y levantar el middleware en su propio equipo siguiendo los comandos del apartado de despliegue y configuración del entorno.

% ============================================================
%  Anexo F — Repositorio del código fuente de ELAN 7.1 modificado
% ============================================================

\section*{\normalfont ANEXO F. Repositorio del código fuente de ELAN 7.1 modificado}
\addcontentsline{toc}{section}{ANEXO F. Repositorio del código fuente de ELAN 7.1 modificado}
\label{anx:repo_elan}

El árbol fuente de ELAN 7.1 con el reconocedor de IA incorporado está disponible en el siguiente repositorio de GitHub:

% TODO: reemplazar por la URL real del repositorio
\begin{center}
\url{https://github.com/TPanchoX/elan-7.1-reconocedor-ia}
\end{center}

Este repositorio parte del código fuente oficial de ELAN 7.1 y agrega el paquete \textbf{mpi.eudico.client.annotator.recognizer.ai}, que reúne el reconocedor, su panel de configuración, el cliente HTTP y los DTO de comunicación con el middleware; las clases originales de la herramienta se mantienen intactas. El procedimiento para compilar este código con Maven, registrar el reconocedor mediante SPI y sustituir el JAR en una instalación existente es el descrito en el Anexo~IV.

% ============================================================
%  Anexo G — Manuales de uso y de desarrollo
% ============================================================

\section*{\normalfont ANEXO G. Manuales de uso y de desarrollo del sistema}
\addcontentsline{toc}{section}{ANEXO G. Manuales de uso y de desarrollo del sistema}
\label{anx:manuales}

Los manuales elaborados como parte de este trabajo se encuentran disponibles en la siguiente carpeta compartida de OneDrive:

% TODO: reemplazar por la URL real de la carpeta compartida
\begin{center}
\url{https://1drv.ms/f/ejemplo-carpeta-manuales-tic}
\end{center}

La carpeta contiene dos documentos. El primero es el manual de usuario dirigido a investigadores, que explica paso a paso cómo poner en marcha el middleware, instalar modelos desde ELAN y ejecutar el reconocedor para obtener anotaciones automáticas. El segundo es el manual técnico dirigido a desarrolladores, que detalla cómo construir y empaquetar nuevos backends de modelos compatibles con el contrato de instalación del sistema. Ambos manuales complementan la información de los Anexos I a IV con un enfoque práctico y con capturas del uso real de la herramienta.
